% BEGIN EMBEDDED BIBLIOGRAPHY
% Bibliography data are part of this source; the compiler writes its temporary input.
\begin{filecontents*}[overwrite]{\jobname.bib}
@unpublished{QED3Manuscript,
  title = {The field-strength {Osterwalder--Schrader} axioms for {QED} in three dimensions},
  date = {2026},
  shorthand = {Cur},
  note = {Comprehensive edition, 263 pp., 20 October 2026, before the modular integration revision}
}
\end{filecontents*}
% END EMBEDDED BIBLIOGRAPHY
% MBM1: full proof companion to the separately maintained MBO1 outline.
\documentclass[11pt,letterpaper,reqno]{amsart}
% BEGIN INLINE manuscript-style
\makeatletter
\NeedsTeXFormat{LaTeX2e}[2020/10/01]
% Use with \documentclass[11pt,letterpaper,reqno]{amsart}.
% This file controls presentation only. It supplies no mathematical content.

\RequirePackage[T1]{fontenc}
\RequirePackage{lmodern}
\RequirePackage{microtype}
\RequirePackage[letterpaper,margin=1in,headheight=15pt]{geometry}
\RequirePackage{amsmath,amsthm,amssymb,mathtools}
\RequirePackage{enumitem}
\RequirePackage{perpage}
\RequirePackage{xcolor}
\definecolor{manuscriptred}{rgb}{0.6,0.1,0.1}
\definecolor{manuscriptblue}{rgb}{0.1,0.3,0.5}
\definecolor{manuscriptmagenta}{HTML}{60445C}
\RequirePackage[breaklinks=true,colorlinks=true,
  linkcolor=manuscriptred,citecolor=manuscriptblue,urlcolor=manuscriptmagenta]{hyperref}
\RequirePackage{cleveref}
\hypersetup{pdfauthor={}}

% Retain the AMS paragraph, title, heading and running-head conventions.
% Equations retain one global Arabic sequence, as in the reference.
\setcounter{tocdepth}{1}
\renewcommand{\thefigure}{\Roman{figure}}
\renewcommand{\thetable}{\Roman{table}}
\MakePerPage{footnote}
\renewcommand{\thefootnote}{%
  \ifnum\value{footnote}>9\arabic{footnote}\else\fnsymbol{footnote}\fi}

% Upright statements. Main results are A, B, C; supporting statements
% share a section-local sequence. These end marks indicate typography,
% never proof completeness, validation, or a deferred-proof status.
\theoremstyle{definition}
\newtheorem{ms@theorem}{Theorem}
\renewcommand{\thems@theorem}{\Alph{ms@theorem}}
\newtheorem*{ms@theoremstar}{Theorem}
\newtheorem{ms@proposition}{Proposition}[section]
\newtheorem{ms@lemma}[ms@proposition]{Lemma}
\newtheorem{ms@definition}[ms@proposition]{Definition}
\newtheorem{ms@corollary}[ms@proposition]{Corollary}
\newtheorem*{ms@corollarystar}{Corollary}

\newcommand{\ms@majorstart}{%
  \pushQED{\qed}\renewcommand{\qedsymbol}{$\blacksquare\square$}}
\newcommand{\ms@minorstart}{%
  \pushQED{\qed}\renewcommand{\qedsymbol}{$\blacksquare$}}
\newenvironment{theorem}[1][]{\ms@majorstart\ms@theorem[#1]}
  {\popQED\endms@theorem}
\newenvironment{theorem*}[1][]{\ms@majorstart\ms@theoremstar[#1]}
  {\popQED\endms@theoremstar}
\newenvironment{proposition}[1][]{\ms@minorstart\ms@proposition[#1]}
  {\popQED\endms@proposition}
\newenvironment{lemma}[1][]{\ms@minorstart\ms@lemma[#1]}
  {\popQED\endms@lemma}
\newenvironment{definition}[1][]{\ms@minorstart\ms@definition[#1]}
  {\popQED\endms@definition}
\newenvironment{corollary}[1][]{\ms@majorstart\ms@corollary[#1]}
  {\popQED\endms@corollary}
\newenvironment{corollary*}[1][]{\ms@majorstart\ms@corollarystar[#1]}
  {\popQED\endms@corollarystar}
% The source's lemma/proposition variants are deliberately not carried over:
% choose one marker for each statement type rather than encode proof status.

\theoremstyle{remark}
\newtheorem{remark}[ms@proposition]{Remark}
\newtheorem*{remark*}{Remark}
% The ordinary amsthm proof environment retains its open square.

\crefname{ms@theorem}{theorem}{theorems}
\Crefname{ms@theorem}{Theorem}{Theorems}
\crefname{ms@proposition}{proposition}{propositions}
\Crefname{ms@proposition}{Proposition}{Propositions}
\crefname{ms@lemma}{lemma}{lemmas}
\Crefname{ms@lemma}{Lemma}{Lemmas}
\crefname{ms@definition}{definition}{definitions}
\Crefname{ms@definition}{Definition}{Definitions}
\crefname{ms@corollary}{corollary}{corollaries}
\Crefname{ms@corollary}{Corollary}{Corollaries}
\crefname{remark}{remark}{remarks}
\Crefname{remark}{Remark}{Remarks}

% Bibliography configuration for this manuscript.
  \RequirePackage[backend=biber,style=alphabetic,backref=true,
    giveninits=true,isbn=false,date=year]{biblatex}
  \addbibresource{\jobname.bib}
  \setcounter{biburlnumpenalty}{100}
  \setcounter{biburllcpenalty}{100}
  \setcounter{biburlucpenalty}{100}
  \renewcommand*{\bibfont}{\footnotesize}
  \renewbibmacro{in:}{}
  \DeclareFieldFormat{pages}{#1}
  % Use the public alphabetic style. No private author-hash overrides;
  % repeated author names are printed in full.
\makeatother
% END INLINE manuscript-style
\usepackage{mathrsfs}
\newcommand{\R}{\mathbb R}
\newcommand{\C}{\mathbb C}
\newcommand{\Z}{\mathbb Z}
\newcommand{\dd}{\mathrm d}
\DeclareMathOperator{\supp}{supp}
\title[Maxwell block minimization and Gaussian fluctuations]
 {Maxwell block minimization and Gaussian fluctuations}
\date{}
\hypersetup{pdftitle={Maxwell block minimization and Gaussian fluctuations},pdfauthor={}}
\begin{document}
\begin{abstract}
We construct an Abelian Maxwell block transformation on periodic
three-dimensional lattices. The constrained Landau minimizer has
exponentially decaying kernels, including first differences and
H\"older first differences, uniformly in the mesh and volume.
Local axial coordinates make the fluctuation form uniformly positive;
its positive inverse square root gives an exact Gaussian block
decomposition. We also prove curvature bounds in local real gauges,
locality of polynomially weakened kernels, explicit field-domain
inclusions, and compatibility of the exact operators with odd
periodic covers.
\end{abstract}
\maketitle
\tableofcontents
\section{Introduction}\label{sec:introduction}
Fix an odd integer \(L\ge5\), a power \(M=L^m\ge L^2\), and
\(0<\alpha<1\). At level \(k\ge0\), set
\(\eta=L^{-k}\) and \(N=\eta^{-1}\).
We use real periodic cochains on
\(T_{\eta,\ell}=\eta\Z^3/\ell\Z^3\), with pairing
\(\eta^3\sum\) on each component. Differences are normalized by
the mesh:
\[
 \partial_\mu f(x)=\eta^{-1}(f(x+\eta e_\mu)-f(x)),\qquad
 (d\lambda)_\mu=\partial_\mu\lambda,\qquad
 (d\mathscr A)_{\mu\nu}
   =\partial_\mu\mathscr A_\nu-\partial_\nu\mathscr A_\mu.
\]
Write \(\delta=d^T\) and \(\Delta_{\mathrm{sc}}=\delta d\)
on scalars. Distances between sites are periodic sup distances.
For bonds use their initial sites, and for coordinate supports use
the minimum distance between their points.
The local regularity norms \(\|B\|_{i,W}\), \(i=1,2,3\), mean
the supremum, first-difference supremum, and \(\alpha\)-H\"older
first-difference supremum on \(W\), respectively. Only contained
bonds and difference cells enter these local norms.
Elsewhere \(\|B\|_2\) denotes the quadratic norm; the domain
section explicitly uses \(\|B\|_i\) for the three global
regularity norms.
H\"older differences compare parallel output bonds at distance
in \((0,1]\), with the input column and derivative direction fixed.

For \(z\in\Z^3/\ell\Z^3\), put
\(\Delta_z=\{c:|c-z|_\infty<1/2\}\).
The partition \(M\)-cubes have side \(M\) and centers in
\(M\Z^3/\ell\Z^3\).
The scalar, bond, and plaquette averages to the unit lattice are
\begin{align}
 (\mathbf Q_k\lambda)(z)
   &=\eta^3\sum_{c\in\Delta_z}\lambda(c),\notag\\
 (\mathbf Q_k\mathscr A)_\mu(z)
   &=\eta^3\sum_{c\in\Delta_z}\sum_{j=0}^{N-1}
          \eta\mathscr A_\mu(c+j\eta e_\mu),\label{eq:photon-block-averages}\\
 (\mathbf Q_k^{(2)}F)_{\mu\nu}(z)
   &=\eta^3\sum_{c\in\Delta_z}\sum_{i,j=0}^{N-1}
          \eta^2F_{\mu\nu}(c+i\eta e_\mu+j\eta e_\nu).\notag
\end{align}
All sums are taken on the periodic covering lattice.
Let \(\mathcal R_k\) be the orthogonal projection onto
\(\Delta_{\mathrm{sc}}(\ker\mathbf Q_k)\) in scalar cochains.
The Landau representative \(H_k^{\mathrm{ph}}A\) minimizes
\(\|d\mathscr A\|_2^2\) under
\begin{equation}\label{eq:photon-Landau-constraints}
 \mathbf Q_k\mathscr A=A,\qquad \mathcal R_k\delta\mathscr A=0.
\end{equation}
Set
\begin{equation}\label{eq:photon-coarse-form}
 \Delta_k=(dH_k^{\mathrm{ph}})^TdH_k^{\mathrm{ph}}.
\end{equation}

For one block step on the unit lattice, let
\(B(y)=\{z:|z-y|_\infty<L/2\}\), \(y\in L\Z^3/\ell\Z^3\).
The scalar average is
\(\mathbf Q\lambda(y)=L^{-3}\sum_{z\in B(y)}\lambda(z)\).
On bonds put
\[
 (\mathbf QA)(y,y+Le_\mu)
   =L^{-4}\sum_{z\in B(y)}A(\Gamma_{z,z+Le_\mu}),\qquad
 (\tau A)(y,z)=\tfrac16\sum_\pi A(\Gamma^\pi(y,z)).
\]
Here \(A(\Gamma)\) is the oriented bond sum; the first paths are
straight, while the latter change coordinates in each of the six
orders. Write \(\mathcal Z=\{Z:\mathbf QZ=0,\ \tau Z=0\}\).
For a coarse field \(A'\) on side \(\ell/L\), set
\(A'_L(b)=L^{-1/2}A'(b/L)\).

\begin{theorem}\label{thm:maxwell-block}
Assume that \(\ell/(LM)\) is an odd integer at least \(13\).
There are constants \(C_P^{\mathrm{ph}}>0\) and
\(0<\gamma_P\le1\), depending only on \(L,\alpha\), for which
the following assertions hold uniformly in \(k,\ell\).

The minimizer \(H_k^{\mathrm{ph}}\) exists uniquely, is real
linear, and satisfies \(\mathbf Q_kH_k^{\mathrm{ph}}=1\).
It preserves constants and closed fields, and \(H_0^{\mathrm{ph}}=1\).
Each of its kernel, first-difference kernel, and short H\"older
first-difference kernel is bounded by
\(C_P^{\mathrm{ph}}e^{-\gamma_Pd(b,b')}\).
Also
\[
 |\Delta_k(b,b')|\le C_P^{\mathrm{ph}}e^{-\gamma_Pd(b,b')}.
\]

There is a real coordinate space \(V_k\) with an orthonormal
basis \(\{\Upsilon\}\), each vector supported in an \(L\)-block
or on a noncentral face bond, and a bijection \(C:V_k\to\mathcal Z\).
The basis choices are translates of one fixed block basis, and
\[
 \dim V_k=(2-2L^{-3})\ell^3,\qquad
 \|z\|_2\le\|Cz\|_2\le\sqrt{5L^3}\|z\|_2.
\]
The form \(\widetilde\Delta_k=C^T\Delta_kC\) is strictly positive,
with upper and lower bounds depending only on \(L\), and
\[
 \big|\widetilde\Delta_k^{-1/2}(\Upsilon,\Upsilon')\big|
 \le C_P^{\mathrm{ph}}e^{-\gamma_Pd(\Upsilon,\Upsilon')}.
\]
Here the distance on the right is between basis supports.

For every real \(A'\), the minimum \(A_k^{\min}(A')\) of
\(\langle A,\Delta_kA\rangle\) under
\(\mathbf QA=A'_L,\ \tau A=0\) exists uniquely. With
\(\mathrm{Sq}_k=C\widetilde\Delta_k^{-1/2}\),
\begin{equation}\label{eq:main-block-split}
 \begin{split}
 &\langle A_k^{\min}(A')+\mathrm{Sq}_k\zeta,\,
   \Delta_k(A_k^{\min}(A')+\mathrm{Sq}_k\zeta)\rangle\\
 &\hspace{35mm}=\|dH_{k+1}^{\mathrm{ph}}A'\|_2^2+\|\zeta\|_2^2.
 \end{split}
\end{equation}
The operator at level \(k+1\) on the right is formed on side
\(\ell/L\) and mesh \(\eta/L\). The harmonic part of \(A'\)
is arbitrary. The probability law on the affine constraint space
with density proportional to
\(\exp(-\langle A,\Delta_kA\rangle/2)\), relative to its Euclidean
volume, is the image of the standard product Gaussian \(\mu_I\)
under \(A=A_k^{\min}(A')+\mathrm{Sq}_k\zeta\).
For every \(q\ge1\),
\[
 \mu_I\{\sup_\Upsilon|\zeta(\Upsilon)|\le q\}
   =e^{-\ell^3\delta\varepsilon^0(q)},\qquad
 0\le\delta\varepsilon^0(q)\le12e^{-q^2/2}.
\]

If \(\square\) is a partition \(M\)-cube and
\(\square^\dagger\) its concentric enlargement of side
\((2L+1)M\), then for every real \(A\) there is an independent
real gauge \(\nu\) on this nonwrapping enlargement such that
\[
 \|H_k^{\mathrm{ph}}A-d\nu\|_{i,\square^\dagger}
       \le C_B\|dA\|_\infty,\qquad i=1,2,3,
\]
where \(C_B\) depends only on \(L,M,\alpha\); the subscripts
\(i\) denote the three regularity norms.

Finally, let \(\ell_2=J\ell_1\), with \(J\) odd and both sides
as above, and use the same translated block basis. Periodic
pullback intertwines \(H_k^{\mathrm{ph}},\Delta_k,C,
\widetilde\Delta_k,(\widetilde\Delta_k+xI)^{-1}\) for \(x\ge0\),
\(\widetilde\Delta_k^{-1/2},\mathrm{Sq}_k\), and
\(A_k^{\min}\) on the two tori.
\end{theorem}

We develop the Abelian block argument of
\cite[Section~10]{QED3Manuscript}. The two gauge choices have
different uses. The Landau representative admits a Fourier
description with constants uniform in the fine mesh. The axial
slice has local coordinates on which the curvature energy is
uniformly coercive. The inverse square root of this coordinate
form gives the Gaussian variables in \eqref{eq:main-block-split}.

The uniformity of the Landau bounds requires cancellation at the
zero Fourier alias before the mesh tends to zero. After that
cancellation, the remaining alias sums converge on a fixed complex
strip. A differentiated summability estimate, using \(\alpha<1\),
permits contour shifting for all three kernels. For the
fluctuation form, exponential conjugation bounds the resolvent
uniformly; its integral gives locality of the positive inverse
square root. These arguments are proved in
\cref{subsec:photon-symbol-regularity,subsec:photon-fluctuation-root}.

The remaining sections provide the photon estimates used in a
nonlinear block integration. Curvature controls the reconstructed
field after independent local real gauge changes. Polynomial
weakening along parameter paths preserves exponential decay and
separates linked components; its exact bounds are stated in
\cref{prop:photon-weakened-bounds}. The field-domain theorem,
\cref{prop:photon-domain-preservation}, states the additional
charge inequalities and the margin for complex fluctuations.
Only the exact operators are asserted to commute with periodic
pullback. No fermion operator or interacting estimate is assumed.
\section{The Landau minimizer and locality}\label{sec:photon-primitives}
\subsection{Coarse averages and the minimizer}\label{subsec:photon-Landau-minimizer}
The averages in \eqref{eq:photon-block-averages} preserve
constants and the mean of a bond field, are onto on bond fields, and
satisfy
\begin{equation}\label{eq:photon-averaging-identities}
 \mathbf Q_kd=d\mathbf Q_k,\qquad
 d\mathbf Q_k=\mathbf Q_k^{(2)}d,\qquad
 \|\mathbf Q_k^{(2)}F\|_2\le\|F\|_2.
\end{equation}
On scalars $\mathbf Q_k^T\lambda$ is constant with value $\lambda(z)$
on $\Delta_z$, and $\mathbf Q_k\mathbf Q_k^T=1$.
The differential identities follow by telescoping along the segments
in \eqref{eq:photon-block-averages}. The plaquette average has total
weight one. Jensen's inequality and summation over coarse sites count
each fine plaquette with total weight $\eta^3$, proving its contraction.
To see surjectivity on bonds, give the fine bond leaving a unit block
in direction $\mu$ the value $N A_\mu(z)$, and give the other bonds
value zero. Every segment of length one starting in $\Delta_z$ crosses
exactly one such boundary bond, and its average is $A_\mu(z)$.

\begin{proposition}\label{prop:photon-Landau-minimizer}
The minimizer exists uniquely. It is real and linear, and
$\mathbf Q_kH_k^{\mathrm{ph}}=1$. It preserves constant fields;
if $dA=0$, then $dH_k^{\mathrm{ph}}A=0$. At $k=0$ it is the identity.
Its characterization is
\begin{equation}\label{eq:photon-minimizer-Euler}
 \mathbf Q_k\mathscr A=A,\qquad
 \delta d\mathscr A\in\operatorname{ran}\mathbf Q_k^T
     \text{ on bonds},\qquad
 \Delta_{\mathrm{sc}}\delta\mathscr A
      \in\operatorname{ran}\mathbf Q_k^T\text{ on scalars}.
\end{equation}
\end{proposition}

\begin{proof}
Surjectivity makes the averaging constraint nonempty. The possible
curvatures form an affine subspace of the finite lattice cochain space,
on which the squared norm has a minimum.
Differentiating along the kernel of $\mathbf Q_k$ shows that a
minimizer satisfies
\[
 \langle\delta d\mathscr A,B\rangle=0;
\]
this is the middle
condition in \eqref{eq:photon-minimizer-Euler}. The parallelogram
identity shows that all minimizers have the same curvature.

For a closed periodic bond field the discrete Hodge decomposition is
$B=c+d\lambda$, where $c$ is its mean. One can see this directly in
Fourier variables: at nonzero momentum, closedness says that the
components of $\widehat B$ are proportional to the forward-difference
symbol, and the proportionality factor is
$|\widehat d|^{-2}\widehat d^*\cdot\widehat B$. At zero momentum
only the constant remains. Exact fields have mean zero. Hence a closed
$B\in\ker\mathbf Q_k$ is $d\lambda$ with
$\lambda\in\ker\mathbf Q_k$: average the decomposition, use
\eqref{eq:photon-averaging-identities}, and subtract the constant value
of $\mathbf Q_k\lambda$. Thus the unconstrained minimizers are
$\mathscr A_1+d(\ker\mathbf Q_k)$.

The scalar Laplacian is injective on $\ker\mathbf Q_k$. Indeed,
$\Delta_{\mathrm{sc}}\lambda=0$ gives $d\lambda=0$, so $\lambda$
is constant and its average forces it to be zero. For
$\lambda\in\ker\mathbf Q_k$,
\[
 \mathcal R_k\delta(\mathscr A_1+d\lambda)
       =\mathcal R_k\delta\mathscr A_1+\Delta_{\mathrm{sc}}\lambda.
\]
Since the first term lies in the range
$\Delta_{\mathrm{sc}}(\ker\mathbf Q_k)$, exactly one such gauge
correction imposes the Landau constraint. Moreover
$\mathcal R_kf=0$ is equivalent to
$\Delta_{\mathrm{sc}}f\perp\ker\mathbf Q_k$; this proves the last
condition of \eqref{eq:photon-minimizer-Euler}.
All these equations are real and linear, so uniqueness proves the
reality and linearity of $H_k^{\mathrm{ph}}$.

Constants satisfy the equations. If $A=c+d\lambda$ is closed, then
$c+d\mathbf Q_k^T\lambda$ is a closed fine field with average $A$,
so the minimum curvature energy is zero. This proves preservation of
closed fields. Finally $\mathbf Q_0=1$ gives the level-zero assertion.
\end{proof}

\subsection{The minimizer symbol and uniform regularity}
\label{subsec:photon-symbol-regularity}

The Fourier calculation must separate the zero alias before the
continuum mesh limit is taken. Set $h=(N-1)/2$ and, for $s\in\C$,
\[
 u(s)=N^{-1}\sum_{j=-h}^he^{isj/N},\qquad
 s_N(s)=2N\sin(s/(2N)),\qquad c(s)=2\sin(s/2).
\]
The identity $u(s)s_N(s)=c(s)$ follows by telescoping the geometric
sum. For $q\in\C^3$ define
\[
 d_\mu(q)=N(e^{iq_\mu/N}-1),\quad
 d_\mu^*(q)=N(e^{-iq_\mu/N}-1),\quad
 D(q)=\sum_\mu s_N(q_\mu)^2,\quad U(q)=\prod_\mu u(q_\mu),
\]
and
\[
 V_\mu(q)=e^{i(1-N^{-1})q_\mu/2}u(q_\mu),\quad
 V_\mu^*(q)=V_\mu(-q),\quad
 \rho_\mu(p)=e^{ip_\mu}-1,\quad\rho_\mu^*(p)=e^{-ip_\mu}-1.
\]
The stars in these formulas denote analytic symbols. At real arguments
they are complex conjugates. The aliases are $q_m=p+2\pi m$,
$m\in\{-h,\ldots,h\}^3$, and
\begin{equation}\label{eq:photon-alias-identities}
 V_\mu(q_m)d_\mu(q_m)=\rho_\mu(p),\qquad
 V_\mu^*(q_m)d_\mu^*(q_m)=\rho_\mu^*(p),\qquad
 V_\mu V_\mu^*=u(q_{m,\mu})^2.
\end{equation}
The scalar averaging symbol is $U(q_m)$; on bonds it is
$U(q_m)V_\mu(q_m)$. Thus the Euler equations
\eqref{eq:photon-minimizer-Euler} can be solved separately in each
alias class.

Write $D_0=D(p)$, $U_0=U(p)$, $u_\mu=u(p_\mu)$, and
\begin{align*}
 X_\mu&=\sum_{m\ne0}\frac{U(q_m)^2u(q_{m,\mu})^2}{D(q_m)},&
 Y&=\sum_{m\ne0}\frac{U(q_m)^2}{D(q_m)^2},\\
 N_\mu&=U_0^2u_\mu^2+D_0X_\mu,&
 N_G&=U_0^2+D_0^2Y,\\
 \mathcal N&=\prod_\mu N_\mu,&
 \Pi&=\mathcal N-\sum_\mu s_N(p_\mu)^2X_\mu
                                      \prod_{\nu\ne\mu}N_\nu.
\end{align*}
The coarse form symbol and the scalar multiplier for the Landau
condition are
\begin{align}
 \Lambda_{\mu\nu}
     &=\delta_{\mu\nu}\frac{D_0}{N_\mu}
       -\frac{\rho_\mu\rho_\nu^*U_0^2\mathcal N}
                    {N_\mu N_\nu\Pi},\notag\\
 \omega_\nu&=\frac{\rho_\nu^*D_0U_0^2\mathcal N}{N_\nu N_G\Pi},
                                                    \label{eq:photon-Lambda-symbol}\\
 h_{\mu\nu}(p,m)
    &=\frac{U(q_m)V_\mu^*(q_m)\Lambda_{\mu\nu}}{D(q_m)}
          +\frac{d_\mu(q_m)U(q_m)\omega_\nu}{D(q_m)^2}
                                      \quad(m\ne0).\notag
\end{align}
The zero-alias formula has its apparent pole cancelled:
\begin{equation}\label{eq:photon-zero-alias}
 h_{\mu\nu}(p,0)=\delta_{\mu\nu}\frac{U_0V_\mu^*(p)}{N_\mu}
       +\frac{d_\mu(p)\rho_\nu^*U_0^3
                 (X_\mu-u_\mu^2D_0Y)\mathcal N}
                  {N_\mu N_\nu N_G\Pi}.
\end{equation}
Indeed, $N_\mu-u_\mu^2N_G=D_0(X_\mu-u_\mu^2D_0Y)$, so
\eqref{eq:photon-zero-alias} equals the preceding formula wherever
$D_0\ne0$.

For real $p\in[-\pi,\pi]^3$, the nonzero aliases have $D(q_m)>0$,
and $X_\mu,Y\ge0$. The positivity of $u_\mu$ gives $N_\mu,N_G>0$.
If $p\ne0$, the identity
\begin{equation}\label{eq:photon-Pi-positive}
 \Pi=\sum_\mu\frac{s_N(p_\mu)^2}{D_0}
                 U_0^2u_\mu^2\prod_{\nu\ne\mu}N_\nu
\end{equation}
shows that $\Pi>0$; at zero $N_\mu=N_G=\Pi=1$.
The algebraic identity
$U_0^2\sum_\mu\rho_\mu\rho_\mu^*/N_\mu=D_0\Pi/\mathcal N$
implies $\sum_\mu\rho_\mu^*\Lambda_{\mu\nu}=0$.
Substitution in the alias equations now gives
\[
 \widehat{(H_k^{\mathrm{ph}}A)_\mu}(q_m)
      =\sum_\nu h_{\mu\nu}(p,m)\widehat A_\nu(p),\qquad
 \widehat{(\Delta_kA)_\mu}(p)
      =\sum_\nu\Lambda_{\mu\nu}(p)\widehat A_\nu(p).
\]
To check the first formula, average the displayed alias expression.
The sums including $m=0$ are $N_\mu/D_0$ and $N_G/D_0^2$.
Its two terms give $\widehat A_\mu$ plus opposite multiples of
$\rho_\mu$, so the average is $A$. Its divergence is
$U(q_m)\widehat\omega(p)/D(q_m)$; applying the scalar Laplacian
gives the scalar range condition, and subtracting its gradient from
$D(q_m)\widehat{H_k^{\mathrm{ph}}A}$ gives
$U(q_m)V_\mu^*(q_m)\widehat\Lambda_\mu(p)$.
These are the other two Euler conditions. At $p=0$ only the zero
alias survives, giving the constant field. Uniqueness of the real
minimizer proves the formula. Pairing its bond Euler equation with
$H_k^{\mathrm{ph}}A'$ and using $\mathbf Q_kH_k^{\mathrm{ph}}=1$
identifies the symbol of $\Delta_k$ with $\Lambda$.

There is a strip of fixed width, independent of $N$, in which these
symbols are analytic. Here is the estimate that provides that width
and the derivative bounds. On
$|\operatorname{Re}p_\mu|<5\pi/4$,
$|\operatorname{Im}p_\mu|<1/4$, put
$a=(8/(5\pi))\sin(5\pi/8)$ and $c_h=\cosh(1/8)$.
The elementary complex sine identities give
\[
 \operatorname{Re}s_N(s)^2\ge
       a^2(\operatorname{Re}s)^2-c_h^2(\operatorname{Im}s)^2,
 \qquad |u(s)|\le c_h
\]
in the required alias range. Consequently, for $m\ne0$,
\begin{equation}\label{eq:photon-alias-estimates}
 \operatorname{Re}D(q_m)\ge|m|_2^2,\qquad
 |u(q_{m,\mu})|\le C\langle m_\mu\rangle^{-1},\qquad
 |d_\mu(q_m)|\le C\langle m_\mu\rangle,
\end{equation}
where $\langle t\rangle=\max(1,|t|)$.
For the first estimate use
$|\operatorname{Re}q_{m,\mu}|\ge(3\pi/4)|m_\mu|$ when
$m_\mu\ne0$; the lower bound before simplification is
$a^2(3\pi/4)^2|m|_2^2-3c_h^2/16\ge|m|_2^2$.
The bound on $u$ follows from $u=c/s_N$, and the bound on $d_\mu$
from its sine expression. Thus $X_\mu,Y,N_\mu,N_G,\Pi$ are
uniformly bounded on the larger rectangle, since their alias sums
are bounded by constant multiples of
$\sum_m\prod_\mu\langle m_\mu\rangle^{-2}$.

On its real cube $u_\mu\ge a$, so
$N_\mu\ge a^8$, $N_G\ge a^6$, and
\eqref{eq:photon-Pi-positive} gives $\Pi\ge a^{24}$.
If an analytic function is bounded by $B$ on that rectangle and is
at least $c>0$ on the slightly smaller real cube, apply Schwarz's
lemma to $t\mapsto f(p'+itp'')-f(p')$ on the disk
$|t|<4B/c$. For $|p''|_\infty<c/(16B)$ its argument stays in the
larger rectangle, and $|f(p)-f(p')|\le c/2$.
Apply this to the three denominators. It gives a fixed width $g_2>0$
on $|\operatorname{Re}p_\mu|<5\pi/4-1/4$ where all remain
bounded away from zero. In particular, the zero-alias formula is
analytic there.

In this smaller strip the alias bounds imply
\begin{equation}\label{eq:photon-h-alias-bound}
 |h_{\mu\nu}(p,m)|\le C\left(
       \langle m_\mu\rangle^{-1}|m|_2^{-2}
          +\langle m_\mu\rangle|m|_2^{-4}\right)
                     \prod_\sigma\langle m_\sigma\rangle^{-1}
                 \quad(m\ne0),
\end{equation}
and $|h(p,0)|\le C$. For every $0\le\alpha<1$ this proves
\begin{equation}\label{eq:photon-alias-Holder-sum}
 \sup_p\sum_m|h_{\mu\nu}(p,m)|
       \max\{1,|d_\lambda(q_m)|,
                      |d_\lambda(q_m)||q_m|_1^\alpha\}\le S_\alpha,
\end{equation}
with $S_\alpha$ independent of $N$. In fact, after multiplication
by $\langle m\rangle_\infty^{1+\alpha}$, the first term in
\eqref{eq:photon-h-alias-bound} is bounded by
\[
 C\langle m_\mu\rangle^{-2}
            \prod_{\sigma\ne\mu}\langle m_\sigma\rangle^{-(3-\alpha)/2},
\]
and the second by $C\prod_\sigma\langle m_\sigma\rangle^{-4/3}$.
Both product sums converge; the first is where $\alpha<1$ is used.

For $x\in\eta\Z^3$, set
$K_{\mu\nu}(p;x)=\sum_me^{iq_m\cdot x}h_{\mu\nu}(p,m)$.
It is $2\pi$-periodic in $p$. To check periodicity, away from $D_0=0$
write $E_\mu=N_\mu/D_0$ and $G=N_G/D_0^2$ as sums over all aliases,
and $r=\sum_\mu\rho_\mu\rho_\mu^*/E_\mu$. Then
\[
 \Lambda_{\mu\nu}=\delta_{\mu\nu}/E_\mu
                  -\rho_\mu\rho_\nu^*/(E_\mu E_\nu r),\qquad
 \omega_\nu=\rho_\nu^*/(E_\nu Gr).
\]
Translation of $p$ by $2\pi$ permutes the alias classes modulo
$2\pi N\Z^3$, and every summand, including $e^{iq_m\cdot x}$,
is periodic modulo this latter lattice. Thus $K$ and $\Lambda$
agree on real overlaps of the rectangles; analyticity extends the
identity and defines them on the full strip. Define the covering
kernel by
$H_{\mu\nu}^\infty(x)=(2\pi)^{-3}\int_{[-\pi,\pi]^3}K_{\mu\nu}(p;x)\dd p$.
Shift the contour to $p+i\gamma_H\operatorname{sgn}(x)$, with
$\gamma_H=g_2/2$. The vertical faces cancel by periodicity, and
\eqref{eq:photon-alias-Holder-sum} gives
\[
 |H^\infty(x)|, |\partial H^\infty(x)|,
 |z|^{-\alpha}|\partial H^\infty(x+z)-\partial H^\infty(x)|
          \le C S_\alpha e^{-\gamma_H|x|_1},
 \qquad x,z\in\eta\Z^3,\quad0<|z|\le1.
\]
The last estimate uses
$|e^{iq\cdot z}-1|\le2.5|z|^\alpha|q|_1^\alpha$ on the shifted
contour. The torus kernel is
$H_k^{\mathrm{ph}}(b,b')=\sum_{n\in\Z^3}H^\infty(c-z'+\ell n)$,
with the corresponding direction indices. Summing the exponential
bound over the lifts proves
\begin{equation}\label{eq:photon-kernel-regularity}
 |H_k^{\mathrm{ph}}(b,b')|, |\partial H_k^{\mathrm{ph}}(b,b')|,
 |\delta_\alpha\partial H_k^{\mathrm{ph}}((b,\tilde b);b')|
                 \le C_H(\alpha)e^{-\gamma_Hd(b,b')}.
\end{equation}
Here $b,\tilde b$ have the same direction and their initial points
are at distance at most one. The same contour argument for the bounded
symbol $\Lambda$ gives
\begin{equation}\label{eq:photon-form-kernel-decay}
 |\Delta_k(b,b')|\le C_\Delta e^{-\gamma_Hd(b,b')},\qquad
 \|\Delta_k\|_{2\to2}\le C_\Delta^{(2)}.
\end{equation}
All three constants are independent of $k$ and $\ell$.
For the operator norm, Parseval's formula reduces the estimate to
$\sup_p\|\Lambda(p)\|$. A uniform bound $A_1$ for its nine entries
gives $\|\Lambda(p)\|\le3A_1$, so one may take
$C_\Delta^{(2)}=3A_1$.

\section{Axial coordinates and the Gaussian block identity}\label{sec:axial}
\subsection{Axial coercivity on a block}\label{subsec:photon-axial-coercivity}

For an interior bond $b=(z,z+e_\nu)$, subtract the two
paths from $y$ to its endpoints and append $b$. For each order this
loop bounds the strips belonging to the directions that come after
$\nu$. Their total number of plaquettes is
$\sum_{\mu\text{ after }\nu}|z_\mu-y_\mu|$. If $\tau A=0$,
averaging over the six orders gives
\begin{equation}\label{eq:axial-interior-bond}
 |A(b)|\le\tfrac12\sum_{\mu\ne\nu}|z_\mu-y_\mu|
           \sup_{p\subset B(y)}|dA(p)|
       \le\tfrac{L-1}{2}\sup_{p\subset B(y)}|dA(p)|.
\end{equation}

For a face between $B(y)$ and $B(y+Le_\nu)$, let $a(z)$ be its face
bond values and let $p_c$ be the maximum curvature on plaquettes in
these two blocks and across their common face. A transverse face step
is one plaquette plus the difference of two interior bond values.
Thus $|a(z)-a(z')|\le L|z-z'|_1p_c$.
The defining average $\mathbf QA$ counts each face bond $L$ times
and $L^3(L-1)$ interior bond occurrences. Therefore its face mean
$\bar a$ satisfies
\[
 |\bar a|\le L|(\mathbf QA)(c)|
                       +\tfrac12(L-1)^2p_c.
\]
The average transverse distance from a face site is at most $L-1$,
so $|a(z)-\bar a|\le L(L-1)p_c$. Combining these gives
\begin{equation}\label{eq:axial-local-bond-bound}
 |A(b)|\le\tfrac32L^2p_c+L|(\mathbf QA)(c)|.
\end{equation}
For interior bonds use \eqref{eq:axial-interior-bond} and omit the
coarse-average term. In particular, on the axial fluctuation space,
\begin{equation}\label{eq:axial-fluctuation-coercivity}
 \mathbf QZ=0,\quad\tau Z=0
     \quad\Longrightarrow\quad
 \|Z\|_\infty\le\tfrac32L^2\|dZ\|_\infty,\qquad
 \|Z\|_2^2\le9L^7\|dZ\|_2^2.
\end{equation}
For the second inequality, square the local bound and replace each
local curvature maximum by its sum of squares. An interior plaquette
belongs to the neighborhoods of at most $3L^3+3L^2\le4L^3$ bonds;
a face plaquette to those of at most $L^2$ bonds. Summation gives
$(9/4)L^4\cdot4L^3=9L^7$.

\subsection{Fluctuation coordinates and the positive square root}
\label{subsec:photon-fluctuation-root}

The fluctuations in one block integration satisfy
$\mathbf Q Z=0$ and $\tau Z=0$. We first choose local coordinates
on this slice. The positive inverse square root of the resulting
quadratic form then gives independent Gaussian variables.

For an $L$-block $B(y)$, let $\mathrm{Bd}(y)$ denote its interior
bond functions, with the counting inner product. The functional
$\tau_y$ has rank $L^3-1$: given values $f(x)$ for $x\ne y$, extend
$f(y)=0$ and apply $\tau_y$ to the interior restriction of $df$.
Every path sum telescopes to $f(x)$. Thus
\[
 \dim\ker\tau_y=3L^2(L-1)-(L^3-1)=2L^3-3L^2+1.
\]
The three coordinate reflections preserve $\ker\tau_0$ and act as
commuting orthogonal involutions. Choose an orthonormal basis
$e_i^0$ of simultaneous eigenvectors and translate it to $e_i^y$.
This choice depends only on $L$.

Each face between neighboring blocks has a central bond, on the
line joining their centers. Let
\[
 V_k=\bigoplus_y\ker\tau_y\ \oplus\!\!\!
       \bigoplus_{b\text{ noncentral face bond}}\R\delta_b,
 \qquad
 \{\Upsilon\}=\{e_i^y\}\cup\{\delta_b\}.
\]
Here $\delta_b$ is the point mass on the positively oriented bond
$b$. These are orthonormal coordinates, and
\begin{equation}\label{eq:photon-fluctuation-dimension}
 \dim V_k=d_V\ell^3,\qquad d_V=2-2L^{-3}.
\end{equation}
The support of a basis vector lies in one $L$-block or consists of
the two ends of a bond.

For $z\in V_k$ define $Cz$ to agree with $z$ off the central bonds,
and set, on the central bond $b_c$ of the face
$c=\langle y,y+Le_\mu\rangle$,
\begin{equation}\label{eq:photon-central-bond-elimination}
 (Cz)(b_c)=-\sum_{\substack{b\text{ on }c\\b\ne b_c}}z(b)
                -\sum_{b\text{ interior}}\frac{n_c(b)}Lz(b).
\end{equation}
The second sum is over the bonds of direction $\mu$ in the two
adjacent blocks; $n_c(b)$ counts the straight averaging paths that
contain $b$. Numbering the bonds on such a line by
$t=0,\ldots,2L-2$ gives
$n_c(b)=\min(t+1,2L-1-t)$. The face bond has multiplicity $L$;
every interior bond has multiplicity at most $L-1$. Consequently
\[
 L^4(\mathbf QZ)(c)=L\sum_{b\text{ on }c}Z(b)
                            +\sum_b n_c(b)Z(b).
\]
Only $Z(b_c)$ is unknown in the equation for this face. Central
bonds occur in no axial path and in no other face equation.
It follows that $C$ is a real linear bijection
\begin{equation}\label{eq:photon-coordinate-bijection}
 C:V_k\longrightarrow
       \{Z:\mathbf QZ=0,\ \tau Z=0\}.
\end{equation}

The following coordinate estimates will be used both for the
covariance and for the small-field domain:
\begin{equation}\label{eq:photon-coordinate-estimates}
 \|Cz\|_\infty\le c_C\sup_\Upsilon|z(\Upsilon)|,
 \qquad
 \|z\|_2\le\|Cz\|_2\le(5L^3)^{1/2}\|z\|_2,
 \qquad
 \sum_b|(C\Upsilon)(b)|\le c_1,
\end{equation}
where $c_C=3L^{9/2}$ and $c_1=14L^{3/2}$. To prove the first
bound, normalize the coordinate supremum to one. An interior value
is at most $(2L^3)^{1/2}$ by Cauchy--Schwarz, and a noncentral face
value is at most one. Formula
\eqref{eq:photon-central-bond-elimination} bounds a central value
by $L^2+2L^3(2L^3)^{1/2}\le3L^{9/2}$.
For the second bound, the new central values are supported away
from $z$. A central row has at most $2L^3$ entries, of modulus at
most one, and an entry contributes to at most two such rows.
Cauchy--Schwarz therefore adds at most $4L^3\|z\|_2^2$.
For a face coordinate, $C\delta_b$ has two entries of modulus one.
For an interior unit vector, the interior $\ell^1$ norm is at most
$(3L^3)^{1/2}$, and the total contribution of the six adjacent
central bonds is at most six times this norm. This proves the
last estimate with $c_1=14L^{3/2}$.
The same formula gives
\begin{equation}\label{eq:photon-coordinate-support}
 d(x,x')\le L\quad
 (x\in\supp C\Upsilon,\ x'\in\supp\Upsilon).
\end{equation}

Coordinate sign changes and translations by $M\Z^3$ permute
$\{\pm\Upsilon\}$. For the interior vectors this follows from the
simultaneous eigenvector choice; for face coordinates it follows
from the preservation of central bonds. Fix a cutoff $p_{0,k}>0$;
its charge-dependent value is chosen in
\cref{subsec:photon-domain-bounds}.
The standard Gaussian
$\mu_I$ on $V_k$ and its normalized restriction $\widehat\mu_k$ to
\[
 \mathrm{Box}_k=
 \{z:\ |z(\Upsilon)|\le p_{0,k}\text{ for every }\Upsilon\}
\]
are invariant under these transformations and under $z\mapsto-z$.
Independence of the coordinates gives
\begin{equation}\label{eq:photon-Gaussian-box-mass}
 \mu_I(\mathrm{Box}_k)=\exp(-\ell^3\delta\varepsilon^0_k),
 \qquad
 \delta\varepsilon^0_k=-d_V\log\Pr(|g|\le p_{0,k})
            \le12\exp(-p_{0,k}^2/2),
\end{equation}
when $p_{0,k}\ge1$, where $g$ is standard normal.
Indeed, the Gaussian tail is at most
\[
 2(2\pi)^{-1/2}p_{0,k}^{-1}\exp(-p_{0,k}^2/2);
\]
at
$p_{0,k}=1$ this is less than $1/2$, and
$-\log(1-t)\le2t$ for $t\le1/2$ suffices.

Set
\[
 \widetilde\Delta_k=C^T\Delta_k C,
 \qquad C_k=\widetilde\Delta_k^{-1}.
\]
The first of these definitions makes sense before invertibility
has been proved. From the axial estimate, the contraction of the
plaquette average, and $\mathbf Q_kH_k^{\mathrm{ph}}=1$, we get
\[
 \|z\|_2^2\le\|Cz\|_2^2
 \le9L^7\|dCz\|_2^2
 =9L^7\|\mathbf Q_k^{(2)}dH_k^{\mathrm{ph}}Cz\|_2^2
 \le9L^7\langle z,\widetilde\Delta_kz\rangle.
\]
The upper bound follows from
\eqref{eq:photon-form-kernel-decay} and
\eqref{eq:photon-coordinate-estimates}. Thus, with
\[
 c_P=\max\{9L^7,\ 5L^3 C_\Delta^{(2)}\},
 \qquad
 C_{\mathrm{Dt}}=c_1^2C_\Delta\exp(2\gamma_HL),
\]
we have
\begin{equation}\label{eq:photon-fluctuation-form-estimates}
 c_P^{-1}\le\widetilde\Delta_k\le c_P,
 \qquad
 |\widetilde\Delta_k(\Upsilon,\Upsilon')|
       \le C_{\mathrm{Dt}}\exp(-\gamma_Hd(\Upsilon,\Upsilon')).
\end{equation}
For the second estimate, expand
$\langle C\Upsilon,\Delta_kC\Upsilon'\rangle$ as a double bond
sum. Every contributing pair satisfies
$d(\Upsilon,\Upsilon')\le d(b,b')+2L$ by
\eqref{eq:photon-coordinate-support}; the two $\ell^1$ bounds
then give the asserted constant.

The distance between supports has a triangle defect at most $L$.
We use the following finite-matrix estimate.

\begin{lemma}\label{lem:positive-root-locality}
Let $I$ be finite, let $d$ be symmetric and nonnegative with
$d(a,a)=0$, and suppose
$d(a,c)\le d(a,b)+d(b,c)+\phi$ for some $\phi\ge0$.
Let $T$ be real symmetric, $T\ge t_0I>0$, and
$|T(a,b)|\le C_Te^{-gd(a,b)}$, where $g>0$.
Choose $S>0$ such that
\[
 S\ge e^{g\phi/2}\sup_a\sum_b(d(a,b)+\phi)e^{-gd(a,b)/2},
 \qquad
 \theta=\min\{g/2,t_0/(2C_TS)\}.
\]
Then, for $x\ge0$,
\[
 |(T+xI)^{-1}(a,c)|\le(t_0/2+x)^{-1}e^{-\theta d(a,c)},
 \qquad
 |T^{-1/2}(a,c)|\le\sqrt{2/t_0}\,e^{-\theta d(a,c)}.
\]
\end{lemma}
\begin{proof}
Fix $c$ and conjugate $T$ by the diagonal matrix
$W(a,a)=\exp(\theta d(a,c))$. The error $E=WTW^{-1}-T$ satisfies
\[
 |E(a,b)|\le\theta C_T\exp(g\phi/2)(d(a,b)+\phi)
                                      \exp(-g d(a,b)/2),
\]
because $|d(a,c)-d(b,c)|\le d(a,b)+\phi$ and
$|\exp(u)-1|\le|u|\exp(|u|)$. The Schur test gives
$\|E\|\le t_0/2$. Therefore
$\|(T+x+E)^{-1}\|\le(t_0/2+x)^{-1}$ for $x\ge0$.
Undoing the conjugation yields
\[
 |(T+x)^{-1}(a,c)|\le(t_0/2+x)^{-1}\exp(-\theta d(a,c)).
\]
The scalar integral
$\pi^{-1}\int_0^\infty x^{-1/2}(t+x)^{-1}\,\dd x=t^{-1/2}$
and finite-dimensional diagonalization now give
\begin{equation}\label{eq:photon-root-integral}
 T^{-1/2}=\frac1\pi\int_0^\infty x^{-1/2}(T+x)^{-1}\,\dd x,
 \qquad
 |T^{-1/2}(a,c)|\le(2/t_0)^{1/2}\exp(-\theta d(a,c)).
\end{equation}

\end{proof}

Apply this to $T=\widetilde\Delta_k$, $\phi=L$,
$t_0=c_P^{-1}$ and $g=\gamma_H$. The basis support count gives
\[
 \sup_\Upsilon\sum_{\Upsilon'}\exp(-g d(\Upsilon,\Upsilon'))
       \le c_{\mathrm{bas}}(g),\qquad
 c_{\mathrm{bas}}(g)=2L^3\exp(2gL)(1+6/g)^3.
\]
For example, assign a basis vector to its block center or to an
end of its face bond. At most $2L^3$ vectors are assigned to any
site, and replacing support distance by point distance costs at
most $2L$; the remaining lattice sum is bounded by $(1+6/g)^3$.
Since $t\exp(-gt/4)\le2/g$, we may take
\[
 S_P=\exp(\gamma_HL/2)(L+2/\gamma_H)
                          c_{\mathrm{bas}}(\gamma_H/4),
 \quad
 \gamma_P=\min\{\gamma_H/2,\ (2c_PC_{\mathrm{Dt}}S_P)^{-1}\}.
\]
With
\[
 C_P^{\mathrm{ph}}=\max\{C_H(\alpha),C_{\mathrm{Dt}},(2c_P)^{1/2},4\exp(2)\},
\]
all the kernels in \eqref{eq:photon-kernel-regularity},
\eqref{eq:photon-fluctuation-form-estimates}, and the positive
square root satisfy their respective bounds at rate $\gamma_P$
with constant $C_P^{\mathrm{ph}}$. In particular,
\begin{equation}\label{eq:photon-square-root-decay}
 C_k^{1/2}=\widetilde\Delta_k^{-1/2},\qquad
 |C_k^{1/2}(\Upsilon,\Upsilon')|
                  \le C_P^{\mathrm{ph}}\exp(-\gamma_Pd(\Upsilon,\Upsilon')).
\end{equation}
At $k=0$, $H_0^{\mathrm{ph}}=1$; its first differences and their
H\"older differences are supported within distance two and have
entries of modulus at most two. Our choice $C_P^{\mathrm{ph}}\ge4\exp(2)$
therefore includes this level as well.

\subsection{The exact block decomposition}
\label{subsec:photon-block-decomposition}

For a real unit bond field $A'$ on the torus of side $\ell/L$,
put $A'_L(b)=L^{-1/2}A'(b/L)$. Let $A_k^{\min}(A')$ minimize
$\langle A,\Delta_kA\rangle$ under
\[
 \mathbf QA=A'_L,\qquad \tau A=0.
\]
This affine space is nonempty: set each face bond of the face $c$
equal to $L A'_L(c)$, and set every interior bond to zero.
Its direction is $C(V_k)$, and
\eqref{eq:photon-fluctuation-form-estimates} makes the quadratic
form positive definite on this direction. The minimizer is
therefore unique and is characterized by
\begin{equation}\label{eq:photon-block-minimizer-orthogonality}
 \langle A_k^{\min}(A'),\Delta_kCz\rangle=0
                    \quad(z\in V_k).
\end{equation}
This characterization is a real linear system with a unique
solution, so $A_k^{\min}$ is real and linear.

Define $\mathrm{Sq}_k=CC_k^{1/2}$. It maps $V_k$ bijectively onto
the axial fluctuations, has zero block average and mean, and
\[
 \langle\mathrm{Sq}_k\zeta,\Delta_k\mathrm{Sq}_k\zeta'\rangle
  =\langle\zeta,C_k^{1/2}\widetilde\Delta_kC_k^{1/2}\zeta'\rangle
  =\langle\zeta,\zeta'\rangle.
\]
Thus the exact Gaussian decomposition is
\begin{equation}\label{eq:photon-Gaussian-block-decomposition}
 \langle A_k^{\min}(A')+\mathrm{Sq}_k\zeta,
           \Delta_k(A_k^{\min}(A')+\mathrm{Sq}_k\zeta)\rangle
       =\langle A_k^{\min}(A'),\Delta_kA_k^{\min}(A')\rangle
                                                   +\|\zeta\|_2^2.
\end{equation}
The density $\exp(-\langle A,\Delta_kA\rangle/2)$ on the affine
constraint space becomes a constant times $\exp(-\|\zeta\|_2^2/2)$.
The affine change of variables has a constant Jacobian, so
normalization gives the standard Gaussian in the $\zeta$ coordinates.
Restricting it to $\mathrm{Box}_k$ retains exactly the mass in
\eqref{eq:photon-Gaussian-box-mass}.

We also need the relation between the two backgrounds. With the
scalar scaling $\lambda_L(x)=L^{1/2}\lambda(x/L)$, there is a
unique real scalar $\omega_k(A')$ with
$\mathbf Q_{k+1}\omega_k(A')=0$ such that
\begin{equation}\label{eq:photon-background-one-step}
 H_k^{\mathrm{ph}}A_k^{\min}(A')
       =(H_{k+1}^{\mathrm{ph}}A')_L-d\omega_k(A'),
 \qquad
 \|dH_k^{\mathrm{ph}}A_k^{\min}(A')\|_2
       =\|dH_{k+1}^{\mathrm{ph}}A'\|_2.
\end{equation}
The harmonic part of $A'$ is arbitrary.
Here $\mathbf Q_{k+1}=\mathbf Q\mathbf Q_k$ is the average from
the mesh-$\eta_k$ lattice of side $\ell$ to its $L$-spaced
lattice. The operator $H_{k+1}^{\mathrm{ph}}$ is formed at mesh
$\eta_{k+1}$ and side $\ell/L$, before scaling by $L$.

Here is the minimizing argument. At mesh $\eta$, the scaled field
$(H_{k+1}^{\mathrm{ph}}A')_L$ minimizes the curvature norm under
the $(k+1)$-step constraint $\mathbf Q_{k+1}\mathscr A=A'_L$.
This follows by scaling the defining problem; the curvature
$L^2$ norm is unchanged. The field
$\mathscr B=H_k^{\mathrm{ph}}A_k^{\min}(A')$ has the same average.
If $\mathscr A$ is any competitor, set $A=\mathbf Q_k\mathscr A$.
Choose the one-step restricted gauge
\[
 \mu(x)=(\tau A)(y,x)
            -L^{-3}\sum_{x'\in B(y)}(\tau A)(y,x'),
                 \qquad x\in B(y).
\]
Then $\mathbf Q\mu=0$ and $\tau(A-d\mu)=0$, by telescoping
the exact-field path sums. Therefore
\[
 \|d\mathscr A\|_2^2
 \ge\|dH_k^{\mathrm{ph}}A\|_2^2
 =\|dH_k^{\mathrm{ph}}(A-d\mu)\|_2^2
 \ge\|d\mathscr B\|_2^2.
\]
Both backgrounds are minimizers of the same problem. The
Landau-minimizer argument shows that their difference is the
coboundary of a scalar with zero $(k+1)$-step average. This scalar
is unique: a closed field with zero average is exact, and the
only constant scalar with zero block average is zero. The two
backgrounds are real and linear in $A'$, hence so is $\omega_k$.
This proves \eqref{eq:photon-background-one-step} and identifies
the first term on the right of
\eqref{eq:photon-Gaussian-block-decomposition} with
$\|dH_{k+1}^{\mathrm{ph}}A'\|_2^2$.

\section{Curvature control in local gauges}
\label{subsec:photon-local-gauge}

Assume from now on that $\ell/(LM)$ is odd and at least $13$.
For an $M$-cube $\square$, let $\square^\dagger$ be its enlargement
by $LM$ in each direction. Its side is $s_d=(2L+1)M\le3\ell/13$.
For a bond field $B$ on a set $W$, write
\[
 \|B\|_{1,W}=\|B\|_{\infty,W},\qquad
 \|B\|_{2,W}=\|\partial B\|_{\infty,W},\qquad
 \|B\|_{3,W}=\|\delta_\alpha\partial B\|_{\infty,W}.
\]
Only bonds and differences contained in $W$ enter these norms.
They are ordinary differences of the photon coefficients; no
connection occurs in their definition.

\begin{proposition}\label{prop:photon-curvature-local-regularity}
For every real unit bond field $A$, including an arbitrary
constant part, and every $M$-cube $\square$, there is a real scalar
$\nu$ on $\square^\dagger$ such that
\begin{equation}\label{eq:photon-curvature-local-regularity}
 \|H_k^{\mathrm{ph}}A-d\nu\|_{i,\square^\dagger}
                         \le C_B\|dA\|_\infty,
                         \qquad i=1,2,3.
\end{equation}
The constant $C_B$ depends on $L,M,\alpha$ and is independent of
$k,\ell$ and the harmonic part of $A$.
\end{proposition}

\begin{proof}
Let $a$ be the mean of $A$, put $A^0=A-a$, and choose the center
$x_0$ of $\square$ as base point. On the canonical lifted period
cell
\[
 W=x_0+\{-(\ell-1)/2,\ldots,(\ell-1)/2\}^3
\]
define $\lambda(x)$ to be the line sum of $A^0$ along the
rectilinear path from $x_0$ to $x$, in coordinate order $1,2,3$.
This defines a scalar on the torus. Put $A''=A^0-d\lambda$ and
$p_A=\|dA\|_\infty$. The path sums of $A''$ vanish. For a bond
whose two ends lie in $W$ without crossing its boundary,
the strip identity used in the axial estimate gives
\begin{equation}\label{eq:photon-inner-basepoint-bound}
 |A''(x,x+e_\mu)|\le\sum_{\nu>\mu}|x_\nu-x_{0,\nu}|p_A
                                      \le2d(x,x_0)p_A.
\end{equation}
Indeed, commuting the final bond through the later portions of
the two base-point paths sweeps one strip in each direction
$\nu>\mu$. Each strip contributes its plaquette sum, and both
remaining path sums are zero.

A seam bond lies on a closed coordinate line with $\ell-1$
inner bonds. Let $\mathrm{hol}(v)$ be the line sum of $A^0$ in
that direction. Translating this line by one transverse step
changes its sum by at most $\ell p_A$, by summing the intervening
$\ell$ plaquettes. Its mean over the transverse two-torus is zero,
since $A^0$ has mean zero. The mean distance in each coordinate
circle is at most $\ell/4$, hence
$|\mathrm{hol}(v)|\le\ell^2p_A/2$. The line sum of $d\lambda$
around a closed line vanishes. Removing the inner bonds using
\eqref{eq:photon-inner-basepoint-bound} therefore gives
\begin{equation}\label{eq:photon-seam-basepoint-bound}
 |A''(b')|\le|\mathrm{hol}(v)|+(\ell-1)\ell p_A
                                  \le\tfrac32\ell^2p_A
\end{equation}
for the seam value. In particular this construction also gives a
global real gauge with
$\|A^0-d\lambda\|_\infty\le2\ell^2\|dA\|_\infty$.

Since $H_k^{\mathrm{ph}}d\lambda$ is closed and has mean zero,
it is exact, say $d\nu'$. On the nonwrapping cube
$\square^\dagger$, let $x$ denote lifted coordinates and set
$\nu=a\cdot x+\nu'$. Then
$H_k^{\mathrm{ph}}A-d\nu=H_k^{\mathrm{ph}}A''$ there.
If $b\subset\square^\dagger$ and $b'$ is inner, its initial
point is at distance at most $d(b,b')+2$ from an end of $b$.
Thus \eqref{eq:photon-inner-basepoint-bound} gives
\[
 |A''(b')|\le(2d(b,b')+4+s_d)p_A.
\]
For a seam bond,
$d(b,b')\ge(\ell-1-s_d)/2$.

Put $g=\gamma_P$ and
$c_{\mathrm{bd}}(g)=3\exp(3g)(1+6/g)^3$, which bounds the sum
of $\exp(-g d(b,b'))$ over unit bonds. Summing any of the three
kernel bounds against the inner bonds gives
\[
 C_P^{\mathrm{ph}}p_A\big[(s_d+4)c_{\mathrm{bd}}(g)
                         +(2/g)c_{\mathrm{bd}}(g/2)\big],
\]
using $2t\exp(-gt)\le(2/g)\exp(-gt/2)$.
For the seam contribution split off half the decay. Since
$s_d\le3\ell/13$,
\[
 \exp(-g d(b,b'))
 \le\exp(g/4)\exp(-5g\ell/26)\exp(-g d(b,b')/2),
 \qquad
 \ell^2\exp(-5g\ell/26)\le15/g^2.
\]
Together with \eqref{eq:photon-seam-basepoint-bound}, this bounds
the seam contribution by
$62C_P^{\mathrm{ph}}p_Ag^{-2}c_{\mathrm{bd}}(g/2)$.
The proposition follows with
\begin{equation}\label{eq:photon-local-regularity-constant}
 C_B=C_P^{\mathrm{ph}}\big[(s_d+4)c_{\mathrm{bd}}(\gamma_P)
 +(2/\gamma_P+62/\gamma_P^2)c_{\mathrm{bd}}(\gamma_P/2)\big].
\end{equation}
\end{proof}

\section{Weakened photon kernels}
\label{subsec:photon-weakened-kernels}

For $a=M$ or $LM$, let $\pi_a$ consist of the closed cubes of
half-side $a$ centered at $a\Z^3/\ell\Z^3$. These are overlapping
parameter cubes, distinct from the partition $M$-cubes.
Two parameter cubes are linked when they contain a common lattice
site; components below refer to chains of such links. For a set
$S$ of unit sites, let $\mathrm{Cub}_a(S)$ be all parameter cubes
meeting $S$.

Choose the shortest rectilinear path $\mathrm{Pth}(p,p')$ in
coordinate order $1,2,3$. Its signed coordinate offsets are
uniquely determined because $\ell$ is odd. A site belongs to at
most $27$ parameter cubes. If $S$ is connected by nearest-neighbor
steps, then $\mathrm{Cub}_a(S)$ is linked: consecutive sites lie
together in a cube, and all cubes containing either site are
linked to this cube. Along a coordinate segment, at most nine
transverse choices of centers occur. Counting the longitudinal
centers for the three segments of $\mathrm{Pth}(p,p')$ gives
\begin{equation}\label{eq:photon-parameter-path-count}
 \#\mathrm{Cub}_a(\mathrm{Pth}(p,p'))
                            \le9(3d(p,p')/a+9).
\end{equation}
An $L$-block, and the ends of a unit bond, meet at most $27$
cubes of either kind.

We smooth only in the row variable of the minimizer. Let
\[
 \vartheta(r)=
 \begin{cases}0&r\le0,\\3r^2-2r^3&0\le r\le1,\\1&r\ge1,
 \end{cases}
 \qquad
 \varphi(t)=\vartheta(2(3/4-|t|)),
 \qquad
 \chi_c(x)=\prod_{\nu=1}^3\varphi((x_\nu-c_\nu)/M).
\]
The coordinates in the last expression are represented with
modulus at most $\ell/2$. Put $\chi_c(b)=\chi_c(x_b)$, where
$x_b$ is the midpoint of $b$. The identity
$\vartheta(r)+\vartheta(1-r)=1$ implies
$\sum_{j\in\Z}\varphi(t-j)=1$, and therefore
$\chi_c\ge0$ and $\sum_c\chi_c=1$ on the torus.
Moreover $\varphi'$ is bounded by three and is Lipschitz with
constant $24$. Tensor products and the averaged derivative
formula for the discrete difference give, for $\tilde b=b+r$
with $|r|\le1$,
\begin{equation}\label{eq:photon-cutoff-differences}
 |\partial_\nu\chi_c(b)|\le3/M,\qquad
 |\chi_c(\tilde b)-\chi_c(b)|\le9|r|/M,\qquad
 |\partial_\nu\chi_c(\tilde b)-\partial_\nu\chi_c(b)|
                                                     \le42|r|/M^2.
\end{equation}
For the last bound the derivative factor contributes $24/M^2$,
and the two remaining factors contribute $2\cdot9/M^2$.
If $\chi_c(b)\ne0$, both ends of $b$ belong to the cube centered
at $c$.

For a unit bond $b'$ with ends $p'_1,p'_2$, define
\[
 P'(c,b')=\mathrm{Cub}_M
                 (\mathrm{Pth}(c,p'_1)\cup\mathrm{Pth}(c,p'_2)),
 \qquad
 w(b,b';s)=\sum_c\chi_c(b)\prod_{a\in P'(c,b')}s_a.
\]
For an interior basis vector set
$\mathrm{St}(e_i^y)=B(y)$ and $\mathrm{Pt}(e_i^y)=\{y\}$;
for $\delta_b$, both sets are the two ends of $b$. Put
\[
 P''(\Upsilon,\Upsilon')=\mathrm{Cub}_{LM}
 \left(\mathrm{St}(\Upsilon)\cup\mathrm{St}(\Upsilon')
 \cup\!\!\bigcup_{\substack{p\in\mathrm{Pt}(\Upsilon)\\
                           p'\in\mathrm{Pt}(\Upsilon')}}
                                  \mathrm{Pth}(p,p')\right).
\]
The weakened kernels are
\begin{equation}\label{eq:photon-weakened-definitions}
 \begin{split}
 H_k^{\mathrm{ph}}(s)(b,b')&=H_k^{\mathrm{ph}}(b,b')w(b,b';s),
                                                  \quad(k\ge1),\\
 C_k^{1/2}(s)(\Upsilon,\Upsilon')
    &=C_k^{1/2}(\Upsilon,\Upsilon')
                         \prod_{a\in P''(\Upsilon,\Upsilon')}s_a,
                                                  \quad(k\ge0).
 \end{split}
\end{equation}
The two families have parameter sets $\pi_M$ and $\pi_{LM}$,
respectively. When both occur in one expression, $s$ denotes
their pair. Each entry is a polynomial, independent of the
connection and real for real parameters. At $s=1$ the exact
operators are recovered. For general complex parameters, the
second family is defined by the displayed kernel; positivity
or the square-root identity is used only at $s=1$.

Fix $0<\alpha_0\le1/4$ and take
\begin{equation}\label{eq:photon-weakening-scale}
 M=L^m\ge L^2,\qquad
 2\alpha_0n_P\log M\le\gamma_PM,\qquad n_P=402.
\end{equation}
Parameters range over the closed polydisc $|s_a|\le M^{\alpha_0}$.
Such an $M$ exists because $(\log M)/M\to0$.
If $\chi_c(\hat b)\ne0$ for a translate of $b$ by distance at
most two, then $d(c,x_b)\le M+1$; applying
\eqref{eq:photon-parameter-path-count} to both paths gives
\[
 \#P'(c,b')\le234+54d(b,b')/M
                              \le n_P(1+d(b,b')/M).
\]
At most $27$ such centers occur: in each coordinate their
representatives lie in an interval of length $3M/2+5<3M$.
The set $P''$ has two endpoint sets and at most four paths.
Since their endpoints are within $L$ of the respective supports,
the same count gives
\[
 \#P''(\Upsilon,\Upsilon')
 \le378+216/M+108d(\Upsilon,\Upsilon')/(LM)
 \le n_P(1+d(\Upsilon,\Upsilon')/M).
\]
For either product, with corresponding distance $D$, we have
\begin{equation}\label{eq:photon-weakening-product-bound}
 \left|\prod_{a\in P}s_a\right|
 \le M^{\alpha_0n_P}\exp(\alpha_0n_P D\log M/M)
 \le M^{\alpha_0n_P}\exp(\gamma_PD/2).
\end{equation}
This is the loss that the exact exponential decay absorbs.

\begin{proposition}\label{prop:photon-weakened-bounds}
Under \eqref{eq:photon-weakening-scale}, put
$\gamma^P=\gamma_P/2$ and $C^P=116C_P^{\mathrm{ph}}M^{\alpha_0n_P}$.
The three kernels of $H_k^{\mathrm{ph}}(s)$ satisfy
\begin{equation}\label{eq:photon-weakened-regularity}
 |H_k^{\mathrm{ph}}(s)(b,b')|,
 |\partial H_k^{\mathrm{ph}}(s)(b,b')|,
 |\delta_\alpha\partial H_k^{\mathrm{ph}}(s)((b,\tilde b);b')|
                                  \le C^P\exp(-\gamma^Pd(b,b')),
\end{equation}
and
\begin{equation}\label{eq:photon-weakened-root-decay}
 |C_k^{1/2}(s)(\Upsilon,\Upsilon')|
                  \le C^P\exp(-\gamma^Pd(\Upsilon,\Upsilon')).
\end{equation}
These bounds hold on the closed polydiscs.
\end{proposition}

\begin{proof}
The root estimate follows at once from
\eqref{eq:photon-square-root-decay} and
\eqref{eq:photon-weakening-product-bound}.
For the minimizer, fix $b'$ and abbreviate
$H(b)=H_k^{\mathrm{ph}}(b,b')$, $w(b)=w(b,b';s)$ and
$b^+=b+\eta e_\nu$. Put
\[
 \mathrm{Pm}=M^{\alpha_0n_P}\exp(\gamma_Pd(b,b')/2).
\]
For $b,b^+,\tilde b,\tilde b^+$, and the intervening translates,
\eqref{eq:photon-cutoff-differences} and the $27$-center count
give
\[
 |w|\le\mathrm{Pm},\quad |\partial_\nu w|\le9\mathrm{Pm},\quad
 |w(\tilde b^+)-w(b^+)|\le27|r|\mathrm{Pm},\quad
 |\partial_\nu w(\tilde b)-\partial_\nu w(b)|
                                             \le14|r|\mathrm{Pm}.
\]
Here $\eta\le1/3$ and $M\ge9$; all translates are within two
of the initial point. A rectilinear path between $b$ and
$\tilde b=b+r$ has at most $3|r|/\eta$ steps, at distance at
most one from $b$. Summing the first-derivative bound for $H$
along this path gives
\[
 |H(\tilde b)-H(b)|
       \le3|r|C_P^{\mathrm{ph}}\exp(\gamma_P)\exp(-\gamma_Pd(b,b')).
\]
Use the discrete product rule
\[
 \partial_\nu(Hw)(b)
             =(\partial_\nu H)(b)w(b^+)+H(b)\partial_\nu w(b).
\]
It bounds the zeroth and first kernels with factors one and ten.
Subtracting this formula at $b$ and $\tilde b$ gives four terms:
the difference of $\partial_\nu H$, the difference of $w(b^+)$,
the difference of $H$, and the difference of $\partial_\nu w$.
Their bounds have factors $1,27,27\exp(\gamma_P),14$, respectively,
times
$|r|^\alpha C_P^{\mathrm{ph}}\mathrm{Pm}\exp(-\gamma_Pd(b,b'))$.
We used $|r|\le|r|^\alpha$ and $\gamma_P\le1$.
Their sum is at most $116$ times this expression, proving
\eqref{eq:photon-weakened-regularity}.
\end{proof}

At level zero the row smoothing is replaced by the diagonal
family
\begin{equation}\label{eq:photon-weakened-level-zero}
 H_0^{\mathrm{ph}}(s)(b,b')
               =\delta_{bb'}\prod_{a\in\pi_M:\,b\subset a}s_a.
\end{equation}
Every product has at most $27$ factors and at least one. The
first difference has two diagonal terms; the H\"older difference
has four, supported within distance two. Thus their entries are
bounded by $4\exp(2\gamma^P)M^{27\alpha_0}
\exp(-\gamma^Pd(b,b'))$, which is at most the right side of
\eqref{eq:photon-weakened-regularity} by $C_P^{\mathrm{ph}}\ge4\exp(2)$.

The weakening also separates components exactly. Suppose that
$s_a=0$ outside a set $\mathscr Y$ of parameter cubes, and let
$Y_1$ be one linked component of their union. A linked path set
$P$ meeting this component either lies entirely in it, or
contains a cube outside $\mathscr Y$ and its parameter product
vanishes. Both $P'$ and $P''$ are linked and contain every
parameter cube meeting the respective column support; $P''$
also contains every cube meeting the row support. In the first
family, every contributing row cube $c$ contains the fine bond
$b$ and belongs to $P'(c,b')$. It follows that a row or column
supported in $Y_1$ depends only on parameters in that component
and has no nonzero entry at an index outside it. A nonzero row
is supported in some component. The diagonal family has the
same properties because all cubes containing $b$ are linked
through an end of $b$.

Finally, coordinate sign changes preserve the ordered shortest
paths, and translations by the relevant center lattice preserve
the parameter cubes. The exact kernels commute with these
transformations: the averaging, axial and Landau constraints
are preserved, and uniqueness identifies the transformed
minimizers and covariance. The basis transforms by a signed
permutation. The parameter products and $\chi_c$ transform
accordingly. Hence $H_k^{\mathrm{ph}}(s)$ is covariant under
coordinate sign changes and $M$ translations, and
$C_k^{1/2}(s)$ under coordinate sign changes and $LM$
translations, with the parameters transformed as well.

\section{Domains for block integration}
\label{subsec:photon-domain-bounds}

We make explicit the domain conditions used with these operators.
Take $0<\varepsilon\le1/200$, $e_k\in(0,1)$, and put
$t_{i,k}=e_k^{-3/4+i\varepsilon}$ for $i=1,2,3$.
The set $R_k$ consists of complex fine bond fields with
$\|B\|_i<t_{i,k}$ in the three norms. A field belongs to
$c\widetilde R_k$, $0<c\le1$, if it has a decomposition
$B=B_0+B_1$, with $B_1\in cR_k$ and $B_0$ real, such that on
each $\square^\dagger$ the field $B_0$ is gauge equivalent by a real scalar
to a real field with norms strictly below $ct_{i,k}$.
The decomposition and gauges need only exist. For
$\widetilde R_k^-$ the local real representative has the bounds
$t_{i,k}$, while the global term $B_1$ has the bounds
$(1-e_k^\varepsilon)t_{i,k}$. Finally,
\[
 \mathrm{Dom}_k=
 \{A\text{ real unit bond field}:\ H_k^{\mathrm{ph}}A
                                      \in\tfrac14\widetilde R_k\}.
\]
These definitions permit every real harmonic part.

Let $C_w\ge C^P$ and $0<\gamma_w\le\gamma^P$, and choose
$c_\gamma$ to bound both sums
\[
 \sup_\Upsilon\sum_{\Upsilon'}\exp(-\gamma_wd(\Upsilon,\Upsilon')),
 \qquad
 \sup_b\sum_{b'}\exp(-\gamma_wd(b,b'))
\]
over basis elements and unit input bonds, respectively; in the
second supremum $b$ is a fine output bond.
The counts above make $c_\gamma$ independent of $k,\ell$.
Fix $p,p_0>0$ and let $p_{0,k}=(-\log e_k)^{p_0}$ and $p_k=(-\log e_k)^p$.
Assume
\begin{equation}\label{eq:photon-domain-small-charge}
 \begin{split}
 p_{0,k}&\le e_k^{-\varepsilon/4},\qquad
 2C_w^2c_\gamma^2c_Ce_k^{\varepsilon/4}\le1,\\
 8C_Bp_k&\le t_{3,k},\qquad
 8e_k^{3/4-7\varepsilon/2}\le1,\qquad
 L^{3/4}\ge8/3.
 \end{split}
\end{equation}
Since the constants are fixed before the charge, these are
conditions on an upper bound $e_0$, valid for every
$0<e_k\le e_0$.

\begin{theorem}\label{prop:photon-domain-preservation}
With the geometry and weakened operators of the preceding sections,
assume \eqref{eq:photon-weakening-scale} and
\eqref{eq:photon-domain-small-charge}. Then every real unit
bond field with $\|dA\|_\infty\le p_k$ belongs to
$\mathrm{Dom}_k$. For every $\zeta\in\mathrm{Box}_k$ and every
pair of parameters in the closed polydiscs, the field
\[
 \mathscr Z_k(s,\zeta)
       =H_k^{\mathrm{ph}}(s)C C_k^{1/2}(s)\zeta
\]
has all three norms bounded by $\frac12e_k^{-\varepsilon/2}$.
If $\ell/L$ is also an admissible side and
$e_{k+1}=L^{1/2}e_k\in(0,1)$, then
\begin{equation}\label{eq:photon-domain-preserved-one-step}
 A_k^{\min}(A')+\mathrm{Sq}_k\zeta\in\mathrm{Dom}_k
 \quad(A'\in\mathrm{Dom}_{k+1},\ \zeta\in\mathrm{Box}_k).
\end{equation}

If $A\in\widetilde R_k^-$ and $t\in\C$, then
$A+t\mathscr Z_k(s,\zeta)\in\widetilde R_k$ whenever
\begin{equation}\label{eq:photon-fluctuation-analytic-radius}
 |t|\le\mathrm{rad}_1:=\tfrac12e_k^{-3/4+9\varepsilon/2}.
\end{equation}
\end{theorem}

\begin{proof}
The first assertion follows from
\eqref{eq:photon-curvature-local-regularity}: its local real
representative has each norm at most
$C_Bp_k\le t_{3,k}/8<t_{i,k}/4$. Use $B_1=0$ in the domain
decomposition.

For the fluctuation bound, the root estimate and the first
summation bound give
\[
 \sup_\Upsilon|(C_k^{1/2}(s)\zeta)(\Upsilon)|
                         \le C_wc_\gamma e_k^{-\varepsilon/4}.
\]
The coordinate estimate for $C$ holds also for complex
coordinates: evaluating its real bound on the signs of the real
row coefficients gives $\sum_\Upsilon|C(b,\Upsilon)|\le c_C$.
Each norm of $H_k^{\mathrm{ph}}(s)v$ is bounded by
$C_wc_\gamma\|v\|_\infty$, by the respective kernel sum.
Consequently all three norms of $\mathscr Z_k(s,\zeta)$ are at most
\[
 C_w^2c_\gamma^2c_Ce_k^{-\varepsilon/4}
                                     \le\tfrac12e_k^{-\varepsilon/2}.
\]
This proves the complex weakened assertion at every level,
including the diagonal family at level zero.

For \eqref{eq:photon-domain-preserved-one-step}, write
$\mathscr A'=H_{k+1}^{\mathrm{ph}}A'$. It is real and belongs to
$\frac14\widetilde R_{k+1}$. On any coarse enlarged cube there
is a real gauge in which $\mathscr A'$ is represented by a real
field $A^c$ with $i$-th norm below $t_{i,k+1}/2$: add the real
global summand to the local representative. For any fine
$M$-cube $\square$, choose a coarse cube $\square'$ with
$\square\subset L\square'$. Then
$\square^\dagger\subset L(\square'^\dagger)$.
The scaling of the three norms is
$L^{-1/2},L^{-3/2},L^{-3/2-\alpha}$, respectively, and
\[
 t_{i,k+1}=L^{-3/8+i\varepsilon/2}t_{i,k}.
\]
Thus $(A^c)_L$ has $i$-th norm strictly below
$\frac12L^{-7/8+3\varepsilon/2}t_{i,k}
\le\frac12L^{-3/4}t_{i,k}\le3t_{i,k}/16$.
By \eqref{eq:photon-background-one-step}, the new background on
$\square^\dagger$ is gauge equivalent to
$(A^c)_L+\mathscr Z_k(1,\zeta)$. The latter fluctuation is real,
and its norms satisfy
\[
 \tfrac12e_k^{-\varepsilon/2}\le t_{3,k}/16\le t_{i,k}/16
\]
by \eqref{eq:photon-domain-small-charge}. The sum has norms
strictly below $t_{i,k}/4$. This is the required local
representative, again with global summand zero.

Finally, for the radius in \eqref{eq:photon-fluctuation-analytic-radius},
the three norms of $t\mathscr Z_k(s,\zeta)$ are at most
$\frac14e_k^{-3/4+4\varepsilon}<e_k^\varepsilon t_{3,k}$.
This fits the margin in the global summand of $\widetilde R_k^-$,
and proves the asserted inclusion. The fluctuation is polynomial
in the two parameter families and linear in $\zeta$.
\end{proof}

\section{Consistency on odd covers}
\label{subsec:photon-cover-consistency}

Let $\ell_2=J\ell_1$ with $J$ odd and with both sides admissible.
Write $\mathrm{pj}^*$ for periodic pullback, also on the coarse
unit lattices and on fluctuation coordinates. The common
reflection eigenbasis in each block makes the last pullback
well defined. Its image is the subspace of $\ell_1$-periodic
fields, and
\[
 \langle\mathrm{pj}^*f,\mathrm{pj}^*h\rangle_{\ell_2}
                              =J^3\langle f,h\rangle_{\ell_1}.
\]
All exact photon operators commute with this pullback:
\begin{equation}\label{eq:photon-exact-cover-consistency}
 O^{(2)}\mathrm{pj}^*=\mathrm{pj}^*O^{(1)},
 \qquad
 O=H_k^{\mathrm{ph}},\Delta_k,C,\widetilde\Delta_k,
                (\widetilde\Delta_k+x)^{-1},C_k^{1/2},\mathrm{Sq}_k.
\end{equation}
The superscript specifies the torus; $x\ge0$. Also
\begin{equation}\label{eq:photon-minimizer-cover-consistency}
 (A_k^{\min})^{(2)}(\mathrm{pj}^*A')
                         =\mathrm{pj}^*(A_k^{\min})^{(1)}(A').
\end{equation}

To prove these identities, first note that $d,\delta$,
$\Delta_{\mathrm{sc}},\mathbf Q_j,\mathbf Q_j^T,\tau,C$
have the same finite row formulas on the periodic lifts. They
commute with pullback and with translations by $\ell_1\Z^3$.
Pulling back the three Euler equations
\eqref{eq:photon-minimizer-Euler} therefore characterizes the
larger-torus minimizer with the pulled-back datum. Uniqueness
gives the identity for $H_k^{\mathrm{ph}}$.

The exact operators intertwine the orthogonal translations by
$\ell_1\Z^3$. Transposing this relation shows that their adjoints
also intertwine these translations and preserve the periodic subspaces. If
$O^{(2)}\mathrm{pj}^*=\mathrm{pj}^*O^{(1)}$, pairing with two
pullbacks and canceling $J^3$ gives the same identity for $O^T$.
Products give the identities for $\Delta_k$ and
$\widetilde\Delta_k$. Invertibility gives the resolvent identity;
integrating it by \eqref{eq:photon-root-integral} gives the root
identity, and hence the identity for $\mathrm{Sq}_k$.

Finally let $A=(A_k^{\min})^{(1)}(A')$ and test its pullback
against a larger-torus axial fluctuation $\widetilde Z$. Sum the
translates of $\widetilde Z$ over
$\ell_1\Z^3/\ell_2\Z^3$ to obtain $\mathrm{pj}^*Z$. Its average
and axial sums vanish, so $Z\in C(V_k)$ on the smaller torus.
The field $\Delta_k^{(2)}\mathrm{pj}^*A$ is periodic; consequently
\[
 \langle\mathrm{pj}^*A,\Delta_k^{(2)}\widetilde Z\rangle_{\ell_2}
  =J^{-3}\langle\mathrm{pj}^*\Delta_k^{(1)}A,
                                         \mathrm{pj}^*Z\rangle_{\ell_2}
  =\langle A,\Delta_k^{(1)}Z\rangle_{\ell_1}=0.
\]
Its average and axial constraints also pull back. The
characterization \eqref{eq:photon-block-minimizer-orthogonality}
proves \eqref{eq:photon-minimizer-cover-consistency}.



\begin{proof}[Proof of \cref{thm:maxwell-block}]
\Cref{prop:photon-Landau-minimizer} gives the minimizer and its
algebraic properties. The alias estimates give
\eqref{eq:photon-kernel-regularity} and
\eqref{eq:photon-form-kernel-decay}. The coordinate construction
gives \eqref{eq:photon-coordinate-bijection},
\eqref{eq:photon-fluctuation-dimension} and
\eqref{eq:photon-coordinate-estimates}; its positive form and root
satisfy \eqref{eq:photon-fluctuation-form-estimates} and
\eqref{eq:photon-square-root-decay}.
Equations \eqref{eq:photon-Gaussian-block-decomposition} and
\eqref{eq:photon-background-one-step} give
\eqref{eq:main-block-split}, with Gaussian restriction mass
\eqref{eq:photon-Gaussian-box-mass}.
\Cref{prop:photon-curvature-local-regularity} gives the local
real gauges. The exact cover identities are
\eqref{eq:photon-exact-cover-consistency} and
\eqref{eq:photon-minimizer-cover-consistency}.
\end{proof}
\printbibliography
\end{document}
